% CV de Vicente Contreras Duarte. Compilar con `pnpm cv` (Tectonic) → public/cv.pdf.
% Pensado para filtros ATS: una columna, sin tablas, íconos ni gráficos,
% títulos estándar y URLs escritas como texto.
\documentclass[10pt,a4paper]{article}

\usepackage{fontspec}
\setmainfont{texgyreheros}[
  Extension=.otf,
  UprightFont=*-regular,
  BoldFont=*-bold,
  ItalicFont=*-italic,
  BoldItalicFont=*-bolditalic
]
\usepackage[spanish]{babel}
\usepackage[margin=1.5cm,top=1.1cm,bottom=1.1cm]{geometry}
\usepackage{enumitem}
\usepackage{titlesec}
\usepackage[hidelinks]{hyperref}

\hypersetup{
  pdftitle={CV - Vicente Contreras Duarte},
  pdfauthor={Vicente Contreras Duarte},
  pdfsubject={Desarrollador de Software | Soporte e Infraestructura TI},
  pdfkeywords={Desarrollo de software, Soporte TI, Infraestructura, Microsoft 365, TypeScript, React, Node.js, Python, SQL, Power BI, ETL}
}

\pagestyle{empty}
\setlength{\parindent}{0pt}
\setlength{\parskip}{0pt}
\linespread{1.02}
\hyphenpenalty=10000
\exhyphenpenalty=10000

\titleformat{\section}{\large\bfseries}{}{0pt}{}[\vspace{2pt}\hrule]
\titlespacing*{\section}{0pt}{8pt}{5pt}

\setlist[itemize]{leftmargin=1.1em, itemsep=1pt, topsep=2pt, parsep=0pt, label=\textbullet}

% \entry{Cargo o título}{Fechas}{Organización · lugar}
\newcommand{\entry}[3]{%
  \textbf{#1}\hfill #2\par
  #3\par
}

\begin{document}
\raggedright

{\LARGE\bfseries Vicente Contreras Duarte}\par
\vspace{3pt}
{\large Desarrollador de Software | Soporte e Infraestructura TI}\par
\vspace{4pt}
Chile \textbar{}
\href{mailto:vicente.contreras.dev@gmail.com}{vicente.contreras.dev@gmail.com} \textbar{}
\href{https://vduarte.dev}{vduarte.dev}\par
\href{https://www.linkedin.com/in/vicente-contreras-duarte-ab7b72343/}{linkedin.com/in/vicente-contreras-duarte-ab7b72343} \textbar{}
\href{https://github.com/VincentPelos}{github.com/VincentPelos}\par

\section{Perfil}
Desarrollador de software y técnico de soporte e infraestructura TI, estudiante de Ingeniería en Informática.
En la Municipalidad de Placilla doy soporte TI, administro Microsoft 365 y desarrollo plataformas internas
y el nuevo sitio web institucional. Tengo experiencia en análisis de datos y Business Intelligence con Python
y Power BI. Abierto a proyectos freelance remotos.

\section{Experiencia}
\entry{Soporte TI}{Mar 2026 -- Actualidad}{Municipalidad de Placilla · Placilla, Chile (presencial)}
\begin{itemize}
  \item Diagnostico y resuelvo incidencias de hardware, software, conectividad, equipos e impresoras.
  \item Configuro redes LAN y Wi-Fi, routers, servidores y servicios institucionales.
  \item Administro y doy soporte a Microsoft 365, SharePoint y Exchange Online.
  \item Participo en la reconstrucción del sitio web institucional y desarrollo herramientas internas de monitoreo y mesa de ayuda.
  \item Automatizo procesos administrativos y técnicos con PowerShell y Node.js.
\end{itemize}
\vspace{4pt}
\entry{Practicante de Soporte TI}{2025}{AIEP · Chile}
\begin{itemize}
  \item Atendí y resolví incidencias de usuarios.
  \item Diagnostiqué y mantuve equipos computacionales; instalé sistemas operativos y software.
  \item Resolví problemas de hardware y conectividad y apoyé tareas de soporte institucional.
\end{itemize}

\section{Proyectos}
\entry{Resiliencia agrícola de Colchagua · Business Intelligence}{2026}{github.com/VincentPelos/colchagua-etl}
\begin{itemize}
  \item ETL en Python (pandas) con datos oficiales (ODEPA, DGA, CR2MET, Banco Central) y Datamart en esquema estrella.
  \item Machine learning y proyecciones validadas fuera de muestra con scikit-learn y statsmodels en Jupyter.
  \item Dashboard en Power BI con DAX: dependencia de la cereza, amenazas climáticas y vulnerabilidad por comuna.
\end{itemize}
\vspace{4pt}
\entry{Sitio web Municipalidad de Placilla · en desarrollo}{2026}{muniplacilla.cl}
\begin{itemize}
  \item Portal institucional responsivo con React, TypeScript, Vite y Tailwind CSS; noticias y documentos en JSON.
  \item Despliegue automatizado con GitHub Actions a cPanel; accesibilidad y normativa del sector público.
\end{itemize}
\vspace{4pt}
\entry{Portafolio vduarte.dev}{2026}{vduarte.dev}
\begin{itemize}
  \item Sitio en Astro y TypeScript con contenido en JSON validado con Zod, animaciones WebGL y Lighthouse 99--100.
  \item CI con GitHub Actions (ESLint, Vitest, build y Docker) y despliegue continuo con Coolify en un VPS propio.
\end{itemize}

\section{Habilidades}
\textbf{Lenguajes:} TypeScript, JavaScript, Python, SQL, PowerShell, HTML, CSS\par
\textbf{Desarrollo web:} React, Node.js, Vite, Tailwind CSS, Astro, APIs REST, SPFx, MySQL\par
\textbf{Datos y BI:} Power BI, DAX, Power Query, pandas, NumPy, scikit-learn, Jupyter, ETL, modelado dimensional\par
\textbf{Infraestructura:} Linux, Windows, Windows Server, TCP/IP, DNS, DHCP, VLAN, SSH, RDP, virtualización\par
\textbf{Microsoft 365:} SharePoint, Exchange Online, Microsoft Entra ID, Microsoft Graph, MSAL, Power Apps\par
\textbf{Herramientas:} Git, GitHub, GitHub Actions, Docker, VS Code, Neovim, WSL2, cPanel\par

\section{Educación y capacitaciones}
\textbf{Ingeniería en Informática, mención Desarrollo de Sistemas} · AIEP\hfill 2025 -- Actualidad\par
\textbf{Programación} (etapa formativa completada) · AIEP\hfill 2023 -- 2025\par

\textbf{Capacitaciones} · Secretaría de Gobierno Digital, Gobierno de Chile (2026): DocDigital (introducción, uso de la plataforma
y gestión avanzada para administradores) y talleres de la plataforma SIMPLE.\par

\end{document}
